% MURL 2026 submission. Compile twice with pdflatex.

\documentclass{article}
\usepackage[nonatbib]{MURL2026}

\usepackage[utf8]{inputenc}
\usepackage[T1]{fontenc}
\usepackage{hyperref}
\usepackage{url}
\usepackage{booktabs}
\usepackage{amsfonts}
\usepackage{amsmath}
\usepackage{amssymb}
\usepackage{nicefrac}
\usepackage{microtype}
\usepackage{xcolor}
\usepackage{graphicx}
\graphicspath{{figures/}}


\title{FairMod: Constrained Reinforcement Learning for\\Equitable Multilingual Content Moderation}

\author{%
  Rishabh Jain \\
  Department of Computer Science \\
  Maynooth University \\
  Maynooth, Co.\ Kildare, Ireland \\
  \texttt{rishabh.jain.2026@mumail.ie}
}


\begin{document}

\maketitle


\begin{abstract}
We frame multilingual Discord content moderation as a Constrained Markov Decision
Process (CMDP). The agent learns a graduated action policy---\textsc{Allow},
\textsc{Warn}, \textsc{Delete}, \textsc{Timeout}, \textsc{Ban}---across 13
languages, with hard action masking enforcing procedural escalation (timeouts
require prior infractions, bans require prior timeouts). The 403-dimensional
observation combines multilingual sentence embeddings, per-language calibrated
XLM-R toxicity scores, user infraction history with cool-down decay, server heat
signals, and a language identity encoding. Fairness is enforced via an adaptive
Lagrangian multiplier penalizing false positives and a Disparate Impact Wrapper
that penalizes any language's ban rate exceeding $1.5\times$ the baseline.
Training is a two-stage behavior-cloning to MaskablePPO pipeline. Trained on
595 synthetic episodes (8,016 steps), the agent achieves 91.7\% accuracy on
\textsc{Delete}/\textsc{Timeout}/\textsc{Ban}---categories where static
baselines score 0--4.5\%---with false-negative rate below $10^{-4}$. Per-language
parity holds for 12 of 14 categories; residual Hindi/Hinglish disparate impact
is traced to upstream XLM-R bias.
\end{abstract}


%=== 1. INTRODUCTION ===========================================================
\section{Introduction}

Online communities face an escalating challenge in managing toxic behavior across
linguistically diverse user bases. Automated content moderation systems traditionally
rely on keyword filters or static toxicity thresholds, approaches that produce coarse
binary decisions (remove or permit) and fail to capture the graduated nature of
community enforcement~[1]. A single ban applied without prior warning violates
procedural norms that users expect of moderation systems; equally, allowing repeated
borderline behavior because no single message clears a fixed threshold enables
sustained harassment campaigns~[2].

Reinforcement learning (RL) offers a compelling alternative: an agent that learns when
to warn, delete, temporarily suspend, and permanently ban based on the full trajectory
of a user's behavior and the real-time toxicity of each message. Prior RL approaches
to moderation have been largely English-centric~[3] and have not addressed the fairness
problem that arises when classifier scores are systematically biased across
languages~[4].

This paper presents a CMDP-based moderation system that addresses these gaps. Our
contributions are fourfold: (i)~a five-action graduated CMDP with hard action masking
that enforces procedural escalation as a constraint of the action space rather than a
soft penalty; (ii)~a 403-dimensional multilingual observation space spanning 13
languages, combining neural embeddings, calibrated toxicity scores, user history with
cool-down decay, and server-level heat signals; (iii)~a dual fairness mechanism (an
adaptive Lagrangian multiplier for false-positive suppression and a Disparate Impact
Wrapper for per-language ban-rate parity); and (iv)~a per-language affine calibration
scheme that corrects systematic XLM-R scoring biases, with a multilingual regex module
providing hard overrides for content the neural classifier under-detects. Source code
is available at \url{https://gitlab.cs.nuim.ie/2026/cs637/proj/fairmod}.


%=== 2. PROBLEM FORMULATION ====================================================
\section{Problem formulation}
\label{sec:formulation}

\subsection{Constrained Markov decision process}

We model the moderation task as a CMDP $\mathcal{M} = (S, A, P, R, \mathbf{C}, \mathbf{d}, \gamma)$,
where $S \subseteq \mathbb{R}^{403}$ is the observation space, $A = \{0,1,2,3,4\}$
is the discrete action space corresponding to \{\textsc{Allow}, \textsc{Warn},
\textsc{Delete}, \textsc{Timeout}, \textsc{Ban}\}, $P$ is the transition kernel,
$R: S \times A \to \mathbb{R}$ is the reward, $\mathbf{C}: S \times A \to \mathbb{R}^K$
is a vector of $K$ constraint cost functions, $\mathbf{d} \in \mathbb{R}^K$ the
constraint thresholds, and $\gamma = 0.95$ the discount factor. We maximize the
return $V^\pi_R = \mathbb{E}_\pi[\sum_t \gamma^t R_t]$ subject to
$V^\pi_{C_k} \leq d_k$ for $k=1,\ldots,K$.

\subsection{Lagrangian relaxation}

We solve the CMDP via Lagrangian relaxation~[5,~14]. The primal--dual objective
$\mathcal{L}(\pi, \boldsymbol{\lambda}) = V^\pi_R - \sum_k \lambda_k (V^\pi_{C_k} - d_k)$
reduces to a standard MDP with augmented per-step reward
$\tilde{R}_t = R_t - \sum_k \lambda_k C_{k,t}$. The dual variable is updated by
projected gradient ascent on the constraint violation:
\begin{equation}\label{eq:lambda_update}
  \lambda_{k+1} = \max\!\left(\epsilon,\;
    \lambda_k + \alpha_\lambda (\hat{V}_{C_k} - d_k)\right),
\end{equation}
with $\hat{V}_{C_k}$ estimated over a 1{,}000-step sliding window,
$\alpha_\lambda = 0.01$, and $\epsilon = 0.01$ preventing $\lambda_k$ collapse.

\subsection{Reward function}

The base reward $R^{\text{base}}_t$ depends on the distance between the chosen
action $a_t$ and a ground-truth best action $a^*_t$ derived from message toxicity
and user history:
\begin{equation}\label{eq:base_reward}
  R^{\text{base}}_t =
  \begin{cases}
    +1.5 & a_t = a^*_t \\
    -0.5 & |a_t - a^*_t| = 1 \\
    -|a_t - a^*_t| & \text{otherwise.}
  \end{cases}
\end{equation}
Context-sensitive penalties augment the base term: $-3.0$ for \textsc{Ban} without
a prior timeout, $-2.0$ for \textsc{Timeout} with $e_t < 2$, $-2.0$ for any punitive
action on benign content ($s_t < 0.20$), and $-3.0$ for \textsc{Allow} on clearly
toxic content ($s_t \geq 0.70$). The combined reward is then scaled and clipped to
$[-1,1]$ via $R_t = \operatorname{clip}(R^{\text{base}}_t/5,\, -1,\, 1)$,
stabilizing the PPO value network~[6].

\subsection{Action masking}

Hard constraints on action legality are implemented via a binary mask
$\mathbf{m}_t \in \{0,1\}^5$, where $m_{t,a} = 0$ renders action $a$ illegal at
step $t$:
\begin{align}
  m_{t,4} &= 0 \;\text{if } h_t[\text{timeouts}] = 0 \label{eq:mask_ban}\\
  m_{t,3} &= 0 \;\text{if } e_t < 2 \label{eq:mask_timeout}\\
  m_{t,a} &= 0 \;\text{if } s_t < 0.15,\; a \in \{1,2,3,4\} \label{eq:mask_low_tox}\\
  m_{t,0} &= 0 \;\text{if } s_t \geq 0.85 \label{eq:mask_high_tox}
\end{align}
where $s_t$ is the calibrated toxicity score, $e_t$ the effective infraction count,
and $h_t[\text{timeouts}]$ the count of prior timeouts. (\ref{eq:mask_low_tox})
blocks all punitive actions \emph{including} \textsc{Warn} when $s_t < 0.15$,
preventing false-positive warnings on clearly benign content. If all actions are
masked, \textsc{Allow} is restored as a safe fallback.


%=== 3. SYSTEM ARCHITECTURE ====================================================
\section{System architecture}
\label{sec:architecture}

The agent receives a 403-dimensional observation
$\mathbf{o}_t = [\mathbf{e}_t \,\|\, s_t \,\|\, \mathbf{h}_t \,\|\, \mathbf{q}_t \,\|\, \mathbf{l}_t]$
at each step (Table~\ref{tab:obs_space}).

\begin{table}[t]
  \caption{Observation space components (403 total dimensions).}
  \label{tab:obs_space}
  \centering
  \begin{tabular}{llcl}
    \toprule
    Symbol & Component & Dims & Description \\
    \midrule
    $\mathbf{e}_t$ & Message embedding    & 384 & MiniLM-L12-v2~[13] \\
    $s_t$          & Toxicity score       & 1   & Calibrated XLM-R score $\in[0,1]$ \\
    $\mathbf{h}_t$ & User history         & 3   & Warns, timeouts, effective infractions \\
    $\mathbf{q}_t$ & Server heat          & 2   & Rolling toxicity rate, action rate \\
    $\mathbf{l}_t$ & Language identity    & 13  & One-hot over 13 language codes \\
    \bottomrule
  \end{tabular}
\end{table}

\paragraph{User infraction model with cool-down decay.}
Each user maintains a ledger tracking cumulative warns $w_u$, timeouts $\tau_u$,
total infractions $n_u$, and a clean-message streak $c_u$. The effective infraction
count decays with sustained clean behavior, implementing rehabilitation:
\begin{equation}\label{eq:decay}
  e_t = \max\!\left(0,\; e_{t-1} - \delta \cdot \max(0,\; c_t - \tau_{\text{cd}})\right),
\end{equation}
where $\tau_{\text{cd}} = 5$ is the clean-streak threshold before decay begins and
$\delta = 1.0$ the decay rate per additional clean message.

\paragraph{Server heat signals.}
Two server-level statistics provide context beyond the individual user: the rolling
mean toxicity score over the last 20 messages, and the fraction of those steps
where a punitive action was taken. These let the agent modulate its response based
on whether the broader channel environment is currently elevated.

\paragraph{Action space.}
Table~\ref{tab:actions} summarizes the five discrete actions and their masking
preconditions.

\begin{table}[t]
  \caption{Action space and masking preconditions.}
  \label{tab:actions}
  \centering
  \begin{tabular}{cll}
    \toprule
    ID & Action & Precondition (must hold to be legal) \\
    \midrule
    0 & \textsc{Allow}   & $s_t < 0.85$ \\
    1 & \textsc{Warn}    & $s_t \geq 0.15$ \\
    2 & \textsc{Delete}  & $s_t \geq 0.15$ \\
    3 & \textsc{Timeout} & $s_t \geq 0.15$,\; $e_t \geq 2$ \\
    4 & \textsc{Ban}     & $s_t \geq 0.15$,\; $h_t[\text{timeouts}] \geq 1$ \\
    \bottomrule
  \end{tabular}
\end{table}


%=== 4. MULTILINGUAL TOXICITY ==================================================
\section{Multilingual toxicity assessment}
\label{sec:toxicity}

\paragraph{Base classifier.}
We use \texttt{textdetox/xlmr-large-toxicity-classifier-v2}~[10] built on
XLM-RoBERTa-large~[9], a 560M-parameter cross-lingual encoder pretrained on
2.5~TB of filtered text across 100 languages. Raw scores
$s_{\text{raw}}(x) \in [0,1]$ are extracted from the TOXIC class probability.
The classifier scores raw message text in isolation---never concatenated context
strings---to prevent lexical contamination from surrounding benign messages.

\paragraph{Per-language affine calibration.}
The XLM-R classifier exhibits systematic per-language score biases: Arabic averages
$0.139$ below English-equivalent content; Italian, Ukrainian, and German are also
under-scored, while Hinglish is over-scored by $+0.095$ on average relative to
English. We address this with per-language affine calibration computed from
percentile anchors in the training corpus. Let $s^l_{\text{low}}$ and
$s^l_{\text{high}}$ be the 75th percentiles of benign and toxic message scores in
language $l$; English anchors serve as the calibration target. The calibrated
score is
\begin{equation}\label{eq:calibration}
  \hat{s}(x, l) = \operatorname{clip}(s_{\text{raw}}(x) \cdot \sigma_l + \beta_l,\; 0,\; 1),
\end{equation}
with $\sigma_l = (s^{\text{en}}_{\text{high}} - s^{\text{en}}_{\text{low}}) /
(s^l_{\text{high}} - s^l_{\text{low}})$ and
$\beta_l = s^{\text{en}}_{\text{low}} - s^l_{\text{low}} \cdot \sigma_l$.
Calibration parameters are computed once from the training corpus and applied
identically during training and production inference.

\paragraph{Multilingual threat detection and gaming-slang override.}
For five threat categories (stalking, violence, death wishes, swatting/doxxing,
mass violence) we maintain language-aware regex patterns across all 13 languages,
with character-based matching for Japanese/Chinese, RTL-safe patterns for
Arabic/Hebrew, and both native-script and Latin transliteration variants for Indic
and Hinglish. When a pattern fires and $\hat{s} < 0.80$, the score is overridden to
$\max(\hat{s}, 0.85)$, ensuring \textsc{Ban} remains available for explicit threats.
Death-threat patterns originally assumed verb-before-object order; adding SOV
reverse-order variants for Russian, Ukrainian, Spanish, and Hindi raised threat
coverage from 9/13 to 13/13. Symmetrically, a compiled regex of common gaming slang
(``gg wp'', ``nice shot'') sets $\hat{s} = 0$ pre-classifier, eliminating a known
class of false positives.


%=== 5. TRAINING METHODOLOGY ===================================================
\section{Training methodology}
\label{sec:training}

\paragraph{Two-stage training: behavior cloning $\to$ RL.}
Recent work on scaling RL for content moderation~[15] demonstrates that RL-Only
training without a behavioral prior is unstable---the policy fails to converge to
a coherent reasoning schema. We therefore adopt a two-stage pipeline. In Stage~1
(behavior cloning), an oracle rule $a^*_t = f(s_t, e_t, \tau_u)$ derived from the
escalation ladder is rolled out across all training episodes; the MaskablePPO
policy is then trained via NLL minimization on the resulting
$(\mathbf{o}_t, a^*_t, \mathbf{m}_t)$ triples for 30 epochs with batch size 256
and learning rate $3 \times 10^{-4}$. In Stage~2 (RL fine-tuning), the
BC-initialized policy is loaded and fine-tuned for $10^6$ MaskablePPO timesteps;
the RL stage can improve \emph{beyond} the oracle---learning, for example, that a
borderline-toxic message from a user with a long clean streak should be treated
more leniently than the rule prescribes.

\paragraph{Base RL algorithm.}
Standard PPO~[6] cannot handle invalid action masking natively. We use
MaskablePPO~[7] from \texttt{stable-baselines3-contrib}~[8], which zeroes out the
log-probabilities of masked actions before the softmax normalization, ensuring
the policy never places probability mass on illegal actions. Hyperparameters:
$n_{\text{steps}} = 4096$, batch size $= 256$, $n_{\text{epochs}} = 10$,
$\gamma = 0.95$, learning rate $= 3 \times 10^{-4}$, entropy coefficient $= 0.01$,
maximum gradient norm $= 0.5$. Four parallel environments yield an effective batch
of $16{,}384$ transitions per policy update, well above the $\sim$1{,}024
inflection point identified by~[15] for stable PPO training on moderation tasks.

\paragraph{Escalation-consistency reward.}
In addition to the per-step accuracy reward (\ref{eq:base_reward}), we add a
\textit{escalation-consistency} component that evaluates trajectory quality~[15]:
\begin{equation}\label{eq:esc_reward}
  R^{\text{esc}}_t =
  \begin{cases}
    +0.3 & a_t > a_{t-1},\; s_t - s_{t-1} > 0.10 \\
    +0.2 & a_t = 0,\; s_t < 0.20,\; c_t \geq 3 \\
    -0.4 & a_t - a_{t-1} \geq 2,\; s_t - s_{t-1} < 0.20 \\
    -0.3 & a_t > 0,\; c_t \geq 5,\; s_t < 0.30 \\
    0    & \text{otherwise.}
  \end{cases}
\end{equation}
This rewards proportional escalation and de-escalation while penalizing
unjustified severity jumps and re-punishment of rehabilitated users.

\paragraph{Wrapper stack.}
The environment is wrapped with four layers, applied in a fixed order so that all
penalty terms accumulate in raw reward space before the final scaling step bounds
the total in $[-1,1]$. The \textbf{LagrangianPenaltyWrapper} implements
(\ref{eq:lambda_update}) with $\lambda_{\text{init}}=0.3$, $\alpha_\lambda=0.01$,
false-positive threshold $d=0.05$ over a 1{,}000-step window (any punitive action
with $s_t < 0.30$ counts as a false positive). The \textbf{DisparateImpactWrapper}
monitors per-language \textsc{Ban} rates over a rolling 2{,}000-step window and
applies a $-2.0$ penalty when language $l$ satisfies
$\hat{\rho}_l > 1.5 \cdot \hat{\rho}_{\text{all}}$ and the current action is
\textsc{Ban}. A \textbf{MetricsTrackingWrapper} records per-episode FP/FN rates
and reward, and the \textbf{RewardScalingWrapper} applies the final clip last.

\paragraph{Train/test split and constraint-aware checkpoint selection.}
Episodes are split 80/20 into train/test; evaluation uses the held-out test split
exclusively. For a CMDP, mean reward alone is the wrong selection criterion---a
model with reward $0.9$ and FP rate $0.08$ is strictly worse than one with reward
$0.85$ and FP rate $0.02$. We therefore use a constraint-aware score
$\bar{R} - 10\cdot\max(0, \overline{\text{FP}} - 0.05)$ to select the best
checkpoint.

\paragraph{Monte-Carlo confidence estimation.}
Single-pass deterministic inference produces bimodal confidence distributions
that overcommit on borderline content~[15]. For calibrated toxicity
$s_t \in [0.30, 0.70]$, we run the policy $N=4$ times under stochastic sampling
and return the majority-vote action $a_t = \operatorname{mode}\{\tilde{a}^{(i)}\}_{i=1}^N$,
breaking ties toward the lower-severity action. Outside this range, deterministic
inference is used.

\paragraph{Synthetic data generation.}
Training episodes are generated in three stages: (i) sourcing a balanced corpus of
$5{,}000$ toxic and $5{,}000$ benign examples per language from the Jigsaw~[11] and
\texttt{textdetox/multilingual\_toxicity\_dataset}~[10] datasets; (ii) $k$-means
clustering on MiniLM embeddings to obtain diverse per-language seeds; and
(iii) prompting DeepSeek~V3.1~[12] via the Baseten API to generate multi-user
Discord threads from those seeds. Thread types follow a fixed mix---30\% toxic,
20\% escalating, 15\% mild, 10\% subtle, 25\% benign---with the Hinglish split
skewed to 20\%/30\% benign/mild via a code-switching system prompt. The final
dataset comprises 595 episodes (8,016 steps) across 14 language categories
(Table~\ref{tab:lang_dist}) under an 80/20 train/test split.

\begin{table}[t]
  \caption{Language distribution in evaluation corpus (8,016 steps).}
  \label{tab:lang_dist}
  \centering
  \begin{tabular}{lrr|lrr}
    \toprule
    Lang & Steps & \% & Lang & Steps & \% \\
    \midrule
    \texttt{hi}  & 1014 & 12.6 & \texttt{uk}  &  478 & 6.0 \\
    \texttt{en}  &  703 &  8.8 & \texttt{de}  &  565 & 7.0 \\
    \texttt{ar}  &  687 &  8.6 & \texttt{zh}  &  502 & 6.3 \\
    \texttt{es}  &  616 &  7.7 & \texttt{ru}  &  484 & 6.0 \\
    \texttt{ja}  &  616 &  7.7 & \texttt{hin} &  400 & 5.0 \\
    \texttt{fr}  &  609 &  7.6 & other        &  247 & 3.1 \\
    \texttt{he}  &  576 &  7.2 & \texttt{it}  &  519 & 6.5 \\
    \bottomrule
  \end{tabular}
\end{table}


%=== 6. EXPERIMENTAL RESULTS ===================================================
\section{Experimental results}
\label{sec:results}

\paragraph{Baseline comparison.}
We compare the RL agent against two baselines on the held-out test split (595
episodes, 8{,}016 steps). The \textbf{keyword filter} triggers a \textsc{Delete}
on hand-crafted slur/profanity matches; all else \textsc{Allow}. The \textbf{static
threshold} maps fixed toxicity ranges to actions: $s < 0.15 \to$ \textsc{Allow},
$0.15 \leq s < 0.50 \to$ \textsc{Warn}, $s \geq 0.50 \to$ \textsc{Delete}, with
no user history considered. Results are in Table~\ref{tab:baseline}. The static
threshold leads on overall accuracy (77.0\%) thanks to perfect \textsc{Allow}
precision, but achieves 0\% on \textsc{Timeout} and \textsc{Ban}---it is
fundamentally incapable of handling recidivism. The RL agent trades 3 points of
overall accuracy for graduated enforcement, with per-class accuracies of 91.6\%
and 89.5\% on \textsc{Timeout} and \textsc{Ban} respectively. Its false-negative
rate ($10^{-4}$) is near zero. The elevated false-positive rate ($0.153$) reflects
conservative over-warning at the benign boundary---the majority of false positives
are \textsc{Warn} actions on benign content, a lower-harm error than false
\textsc{Delete} or \textsc{Ban}.

\begin{table}[t]
  \caption{Baseline performance comparison (8,016 steps, held-out test split).}
  \label{tab:baseline}
  \centering
  \begin{tabular}{lcccc}
    \toprule
    System & Accuracy & FP rate & FN rate & \textsc{Timeout}/\textsc{Ban} \\
    \midrule
    Keyword filter   & 0.673 & 0.020 & 0.272 & none \\
    Static threshold & 0.770 & 0.000 & 0.000 & none \\
    \textbf{RL agent}& \textbf{0.736} & 0.153 & $\mathbf{10^{-4}}$ & \textbf{yes} \\
    \bottomrule
  \end{tabular}
\end{table}

\paragraph{Per-class accuracy.}
Table~\ref{tab:confusion} shows the agent's per-class accuracy on the held-out test
split. \textsc{Delete} (94.1\%), \textsc{Timeout} (91.6\%), and \textsc{Ban}
(89.5\%) all exceed 89\%, demonstrating reliable graduated enforcement.
\textsc{Allow} accuracy (77.8\%) is reduced by conservative over-warning: 1{,}223
of 5{,}530 benign messages receive a \textsc{Warn}. \textsc{Warn} accuracy
(19.4\%) is the weakest class because the agent prefers \textsc{Delete} for most
first offenses, reflecting asymmetric reward penalties where under-escalation is
penalized less than over-escalation.

\begin{table}[t]
  \caption{RL agent per-class accuracy (held-out test split).}
  \label{tab:confusion}
  \centering
  \begin{tabular}{lcc}
    \toprule
    True action & Correct & Per-class accuracy \\
    \midrule
    \textsc{Allow}   & 4305 / 5530 & 77.8\% \\
    \textsc{Warn}    &  185 /  953 & 19.4\% \\
    \textsc{Delete}  &  507 /  539 & 94.1\% \\
    \textsc{Timeout} &  406 /  443 & 91.6\% \\
    \textsc{Ban}     &  493 /  551 & 89.5\% \\
    \midrule
    \textbf{Overall} & \textbf{5896 / 8016} & \textbf{73.6\%} \\
    \bottomrule
  \end{tabular}
\end{table}

\paragraph{Procedural scenario testing.}
Three designed scenarios verify the learned policy's behavioral contracts.
\textbf{Escalation:} 10 messages of monotonically increasing toxicity produce the
full graduated ladder \textsc{Allow}$^3 \to$ \textsc{Warn} $\to$ \textsc{Delete}
$\to$ \textsc{Timeout}$^3 \to$ \textsc{Ban}, respecting all masking preconditions.
\textbf{Sustained troll:} a user alternating borderline-toxic and benign messages
over 10 steps---never triggering a single high-severity event---is correctly
escalated through recidivism (\textsc{Delete} $\to$ \textsc{Allow} $\to$
\textsc{Delete} $\to$ \textsc{Timeout}$^3$), confirming the policy uses cumulative
history rather than per-message severity alone. \textbf{Rehabilitation:} after an
initial \textsc{Delete}, 12 clean messages drive $e_t \to 0$ by step 7, and a
subsequent mild gripe at step 13 produces \textsc{Allow} rather than re-punishment.

\paragraph{Cross-lingual parity.}
Table~\ref{tab:parity} reports the XLM-R cross-lingual spread for six representative
sentence categories. Three categories show spreads $> 0.30$: mild frustration
(Arabic, Italian, Ukrainian, German under-detect), veiled threats (Ukrainian,
Italian score near zero while Japanese scores near one), and severe toxic content
(Arabic scores $0.349$ vs.\ near-unity elsewhere). These reflect fundamental
limitations of the XLM-R pretrained weights rather than calibration failures, as
affine calibration cannot supply signal that the classifier does not produce.

\begin{table}[t]
  \caption{Cross-lingual toxicity score spread after calibration.}
  \label{tab:parity}
  \centering
  \begin{tabular}{llcc}
    \toprule
    Category & Severity & Spread & Status \\
    \midrule
    \texttt{benign\_greeting}   & benign   & 0.001 & \checkmark \\
    \texttt{direct\_insult}     & moderate & 0.067 & \checkmark \\
    \texttt{hate\_speech}       & extreme  & 0.045 & \checkmark \\
    \texttt{severe\_toxic}      & extreme  & 0.650 & $\times$   \\
    \texttt{mild\_frustration}  & mild     & 0.947 & $\times$   \\
    \texttt{threat\_veiled}     & threat   & 0.976 & $\times$   \\
    \bottomrule
  \end{tabular}
\end{table}

\paragraph{Fairness audit.}
Table~\ref{tab:fairness} reports per-language \textsc{Ban} rates (overall
$0.066$). Twelve of fourteen categories pass the $1.5\times$ parity threshold;
Hindi ($1.92\times$) and Hinglish ($1.64\times$) exhibit residual disparate
impact, traced via Table~\ref{tab:parity} to upstream XLM-R bias ($+0.095$ for
Hindi vs.\ English) that compounds through the escalation ladder. Aggressive
calibration offsets did not reduce the disparity, confirming the bias is
structural; classifier fine-tuning is the recommended remediation~[16].

\begin{table}[!h]
  \caption{Per-language \textsc{Ban} rate fairness audit (held-out test split).}
  \label{tab:fairness}
  \centering
  \small
  \begin{tabular}{lccc|lccc}
    \toprule
    Lang & Rate & Ratio & St. & Lang & Rate & Ratio & St. \\
    \midrule
    \texttt{ar}  & 0.074 & 1.13 & \checkmark & \texttt{it}  & 0.004 & 0.06 & \checkmark \\
    \texttt{de}  & 0.050 & 0.76 & \checkmark & \texttt{ja}  & 0.068 & 1.04 & \checkmark \\
    \texttt{en}  & 0.048 & 0.74 & \checkmark & \texttt{ru}  & 0.072 & 1.10 & \checkmark \\
    \texttt{es}  & 0.036 & 0.54 & \checkmark & \texttt{uk}  & 0.082 & 1.24 & \checkmark \\
    \texttt{fr}  & 0.007 & 0.10 & \checkmark & \texttt{zh}  & 0.098 & 1.49 & \checkmark \\
    \texttt{he}  & 0.078 & 1.19 & \checkmark & \texttt{hi}  & 0.126 & 1.92 & $\times$ \\
    \texttt{hin} & 0.108 & 1.64 & $\times$ &              &       &      &  \\
    \bottomrule
  \end{tabular}
\end{table}


%=== 7. CONCLUSION =============================================================
\section{Conclusion}
\label{sec:conclusion}

We presented a CMDP-based multilingual moderation system across 13 languages
with dual fairness constraints. Action masking embeds procedural constraints into
the learned policy without reward engineering, and two-stage BC$\to$RL
training~[15] markedly improves convergence over RL-Only; the RL agent achieves
89--94\% accuracy on \textsc{Delete}/\textsc{Timeout}/\textsc{Ban} actions, where
static baselines score $\leq 4.5\%$. Lagrangian methods plus the
DisparateImpactWrapper constrain fairness for 12 of 14 categories; residual
Hindi/Hinglish disparate impact traces to upstream XLM-R bias that calibration
cannot fully correct. Limitations include the small training corpus (595
episodes) and PPO seed variance (73--82\%); future work covers classifier
fine-tuning, disagreement-based reweighting~[15], and multi-channel server
graph extensions.


\begin{ack}
The author thanks the MURL 2026 course staff for the project framework and
guidance.
\end{ack}


\section*{References}

\medskip
{\small

[1] Watanabe, H., Bouazizi, M.\ \& Ohtsuki, T.\ (2018) Hate speech on Twitter:
A pragmatic approach to collect hateful and offensive expressions and perform
hate speech detection. {\it IEEE Access}, 6:13{,}440--13{,}438.

[2] Habib, R.\ \& Louafi, M.\ (2019) Community membership and the role of
moderators in online platforms. In {\it Proc.\ ACL Workshop on Online Abuse
and Harms}, Florence, Italy.

[3] Gao, C., Chen, J., Liu, S., Lan, Y., Xu, Z.\ \& Cheng, Y.\ (2018)
Reinforcement learning from imperfect demonstrations. In {\it Proc.\ ICLR}.

[4] Risch, I., Leiber, A.\ \& Krestel, R.\ (2020) Bagging BERT models for
robust aggression identification. In {\it Proc.\ Workshop on Trolling,
Aggression and Cyberbullying}.

[5] Altman, E.\ (1999) {\it Constrained Markov Decision Processes}.
CRC Press.

[6] Schulman, J., Wolski, F., Dhariwal, P., Radford, A.\ \& Klimov, O.\ (2017)
Proximal policy optimization algorithms. {\it arXiv preprint arXiv:1707.06347}.

[7] Huang, S.\ \& Onta\~n\'on, S.\ (2022) A closer look at invalid action
masking in policy gradient algorithms. In {\it Proc.\ FLAIRS}.

[8] Raffin, A., Hill, A., Gleave, A., Kanervisto, A., Ernestus, M.\ \& Dormann,
N.\ (2021) Stable-baselines3: Reliable reinforcement learning implementations.
{\it J.\ Mach.\ Learn.\ Res.}, 22(268):1--8.

[9] Conneau, A.\ et al.\ (2020) Unsupervised cross-lingual representation
learning at scale. In {\it Proc.\ ACL}, pp.\ 8440--8451.

[10] Dementieva, D.\ et al.\ (2023) Overview of the 1st shared task on
multilingual text detoxification. In {\it Proc.\ ACL Workshop on NLP for
Positive Impact}.

[11] Google Jigsaw (2018) Toxic comment classification challenge. Kaggle
Competition Dataset.
\url{https://www.kaggle.com/c/jigsaw-toxic-comment-classification-challenge}

[12] DeepSeek-AI (2024) DeepSeek-V3 technical report.
{\it arXiv preprint arXiv:2412.19437}.

[13] Reimers, N.\ \& Gurevych, I.\ (2019) Sentence-BERT: Sentence embeddings
using Siamese BERT-networks. In {\it Proc.\ EMNLP}, pp.\ 3982--3992.

[14] Stooke, A., Achiam, J.\ \& Abbeel, P.\ (2020) Responsive safety in
reinforcement learning by PID Lagrangian methods. In {\it Proc.\ ICML},
pp.\ 9133--9143.

[15] Liu, Z.\ et al.\ (2025) Scaling reinforcement learning for content
moderation with large language models. {\it arXiv preprint arXiv:2512.20061}.

[16] Feldman, M., Friedler, S.~A., Moeller, J., Scheidegger, C.\ \&
Venkatasubramanian, S.\ (2015) Certifying and removing disparate impact.
In {\it Proc.\ ACM SIGKDD}, pp.\ 259--268.

}

\end{document}
